\documentclass[10pt]{extarticle}

\usepackage[a4paper,margin=1cm,left=1cm,right=1cm]{geometry} 
\usepackage[UKenglish]{babel}

\usepackage{amsmath}
\usepackage{amssymb}
\usepackage{amsthm}

\newtheorem{theorem}{Theorem}
\newtheorem{lemma}{Lemma}
\newtheorem{corollary}{Corollary}[theorem]
\newtheorem{proposition}{Proposition}
\theoremstyle{definition}
\newtheorem{definition}{Definition}
\newtheorem{example}{Example}
\theoremstyle{remark}
\newtheorem{remark}{Remark}

\newcommand{\R}{\mathbb{R}}
\newcommand{\Z}{\mathbb{Z}}
\newcommand{\N}{\mathbb{N}}
\newcommand{\Q}{\mathbb{Q}}
\newcommand{\pf}{\mathbb{P}}
\newcommand{\E}{\mathbb{E}}

\DeclareMathOperator{\Cor}{Cor}
\DeclareMathOperator{\Var}{Var}
\DeclareMathOperator{\Std}{Std}
\DeclareMathOperator{\Cov}{Cov}

\usepackage{enumitem}
\usepackage{multicol}
\usepackage{graphicx}
\graphicspath{{./images/}}

\AtBeginDocument{
  \setlength{\abovedisplayskip}{2pt plus 2pt minus 2pt}
  \setlength{\belowdisplayskip}{2pt plus 2pt minus 2pt}
  \setlength{\abovedisplayshortskip}{2pt plus 2pt}
  \setlength{\belowdisplayshortskip}{2pt plus 2pt minus 2pt}
}


\begin{document}

\begin{multicols}{2}

% Probability

    \noindent
    {\large\textbf{Probability}}

    $\Omega$ contains all possible outcomes for random phenomena. 
    $\pf(\cdot)$ maps every event subest $E$ from the sample space to 
    number on $[0, 1]$. This must satisfy: 
    \begin{enumerate}
        \item $\pf(E) \geq 0, \forall E \subset \Omega$,
        \item $\pf(\Omega) = 1$,
        \item $\pf(E_1 \cup E_2) = \pf(E_1) + \pf(E_2) \forall E_1,E_2
            \subset \Omega$ if $\E_1 \cap E_2 = \emptyset$.
    \end{enumerate}

    \noindent
    \textbf{PMF/PDF}
    
    Let $X$ be a DRV. The pmf for $X$ for all $x \in \R$ is \[ 
        f_X(x) = 
        \begin{cases} 
            \pf(X = x), & x \in \Omega, \\
            0, & \text{Otherwise.}
        \end{cases}
    \]
    Valid iff $0 \leq f_X(x) \leq 1, \forall x \in \R;$ and $\sum_{x \in
    \Omega} f_X(x) = 1.$
    
    Let $Y$ be a CRV. The pdf for $Y$ for all $y \in \R$ is such that for any 
    $a, b \in \R$ \[
        \pf(a \leq Y \leq b) = \int_a^b f_Y(y) dy.
    \]
    Note that $\pf(Y=y)=0$.
    Valid iff $f_Y(y) \geq 0, \forall y \in \R;$ and
    $\int_{-\infty}^{\infty}f_Y(y)dy = 1.$

    \noindent
    \textbf{CDF/Quantile}

    The cdf for X (DRV) and Y (CRV) is
    $$F_X(x) = \pf(X \leq x) = \sum_{u \leq x} f_X(u),$$ and 
    $$F_Y(y) = \pf(Y \leq y) = \int_{- \infty}^y f_Y(u) du.$$

    The cdf is non-decreasing func with range $[0,1]$ satisfying 
    $$\lim_{u \to -\infty} F(u) = 0; \lim_{u \to \infty} F(u) = 1;$$ and 
    $$F(u) - F(v) \geq 0, \forall u, v \in \R, u > v.$$

    By FundTh'mCalc have $f_Y(y) = \frac{d}{dy}F_Y(y).$

    Quantile function is inverse of cdf $F_Y(y)$ (valid since it is bjiective),
    $q_p = F_Y^{-1}(p)$ for $p \in [0,1]$.

    \noindent
    \textbf{(Bi/Multi)variate Distributions}

    If $X, Y$ are DRV, then joint pmf is \[ 
        f_{XY}(x,y) = \pf(X = x \cap Y = y).
    \]
    If $X,Y$ are CRV, then joint pdf is def such that they take some value
    within some region $T \subset \R^2$ given by: \[
        \pf((X,Y) \in T) = \iint_{(x,y) \in T} f_{XY}(x,y)dxdy.
    \]

    \textit{Marginal pmf/pdf} describes behaviour of one RV irrespective of 
    potential values for other RV; \[
        f_X(x) = 
        \begin{cases}
            \sum_y f_{XY}(x, y), & \text{ for DRV $Y$,} \\
            \int_{\infty}^{\infty} f_{XY}(x, y)dy, & \text{ for DRV $Y$,}
        \end{cases}
    \]

    \textit{Conditional pmf/pdf} describes how one RV is distributed when
    another takes a particular value. Given $Y = y$: \[
        f_{X | Y} (x | y) = f_X(x) f_Y(y), \text{ where $f_Y(y) > 0$.}
    \]

    \noindent
    \textbf{IID}

    $X,Y$ are independent iff the joint pmf/pdf can be expressed as a product
    of their marginal distributions: \[
        f_{XY}(x,y) = f_X(x)f_Y(y), \; \forall x,y \in \R.
    \]
    In general, \[
        f_\mathbf{X}(\mathbf{x}) = \prod_{i = 1}^{n} f_{X_i}(x_i), \; \forall
        \mathbf{x} \in \R^n.
    \]

    If indep, then $\Cor(X, Y) = 0$. Reverse not nec true. 

    Identically distributed iff they have the same marginal pmf/pdf. 

    If $\mathbf{X}$ is a seq of RV that are iid, then joint pmf/pdf is: 
    \begin{align*}
        f_{\mathbf{X}}(\mathbf{x}) &= f(x_1) \cdots f(x_n) \\
        &=\prod_{i = 1}^{n} f(x_i), \text{ for $\mathbf{x} \in \R^n$}
    \end{align*}

    \noindent
    \textbf{Expectation}
    \[
        \E[X] = \begin{cases}
            \sum_x xf_X(x), & \text{for DRV $X$,} \\
            \int_{-\infty}^{\infty} xf_X(x)dx, & \text{for CRV $X$.}
        \end{cases}
    \]
    Expectation for some transformation of a RV $Z = g(X)$: 
    \[ 
        \E[Z] = \E[g(X)] = \begin{cases}
            \sum_x g(x)f_X(x), & \\
            \int_{-\infty}^{\infty} g(x)f_X(x)dx. & 
        \end{cases}
    \]

    Expectations are invariant to linear transformations: $\E[aX+b] = a\E[X] +
    b$ for constants $a, b \in \R$.

    Expectation of a sum of RVs is the sum of their expectations: 
    $\E\left[ \sum_{i=1}^{n} X_i \right] = \sum_{i=1}^n \E[X_i]$.

    If $X,Y$ are indep RVs, then $\E[XY] = \E[X] \E[Y]$.

    \noindent
    \textbf{Variance}
    \begin{align*}
        \Var(X) &= \E\left[ (X - \E[X])^2\right] \\
                &= \E[X^2] - \E[X]^2 \\
                &= \E[X(X - \E[X])].
    \end{align*}
    Always non-negative. 

    $\Std(X) = \sqrt{\Var(x)}.$

    $\Var(aX + b) = a^2\Var(X), \forall a, b \in \R$.

    If $X_1, \ldots, X_n$ are all independent RV, then: 
    \[ 
        \Var\left(\sum_{i=1}^n X_i\right) = \sum_{i=1}^n \Var(X_i)
    \]

    \noindent
    \textbf{Covariance/Correlation}
    \begin{align*}
        \Cov(X, Y) &= \E[(X - \E[X])(Y - \E[Y])] \\
                   &= \E[XY] - \E[X] \E[Y] \\
        \Cor(X, Y) &= \frac{\Cov(X, Y)}{\sqrt{\Var(X)\Var(Y)}}.
    \end{align*}

    $\Cov(X, X) = \Var(X).$

    When $X, Y$ are indep, then $Cov(X, Y) = 0$. Converse not nec true. 

    $\Cov(aX + b, cY + d) = ac\Cov(X, Y)$, $\forall a, b, c, d \in \R.$

    The variance of a linear combination of two, potentially dependent, RVs is: 
    $\Var(aX + bY + c) = a^2\Var(X) + b^2\Var(Y) + 2ab\Cov(X, Y),$
    for constants $a, b, c \in \R$.

    Correlation is unit-less in the range $[-1,1]$ with the sign indication the
    direction of the linear dependence. Being uncorrelated does not imply
    independence. 

    \noindent
    \textbf{Central limit theorem}

    Let $X_1, \ldots, X_n$ be iid RVs from any distribution with expectation $\mu$ 
    and finite variance $\sigma^2 < \infty$. Then: \[
        \pf(X_1 + \ldots + X_n \leq x) \to \Phi \left(\frac{x -
        n\mu}{\sqrt{n}\sigma}\right),
    \] as $n \to \infty,$ where $\Phi(\cdot)$ denotes the std normal cdf. 
    Therefore, for a sufficiently large sample size $n$: \[ 
        \sum_{i=1}^n X_i \sim \text{N}(n\mu, n\sigma^2), 
    \] or equivalently, \[
        \bar{X} = \frac{1}{n} \sum_{i = 1}^{n} X_i \sim \text{N}\left(\mu,
        \frac{\sigma^2}{n} \right).
    \]

    Generally need $n > 30$ in order to apply the CLT. 

% Basic definitions

    \noindent
    {\large\textbf{Basic definitions}}

    \textit{Population}: the collection of all items that is the focus of what
    is being investigated. 

    \textit{Sample}: any subset of the population from which the data (also
    known as observations and realisations) are collected. 

    \textit{Census}: sample that contains every member of the population
    exactly once. 

    Sample from a population said to be \textit{representative} if all members
    of a population have equal chance of being a member of the sample. 

    \textit{Repeated sampling}: assumptive process of taking multiple samples
    from a population that has same characteristics as observed sample for the
    purpose of assessing uncertainty in these characteristics.

% Estimators

    \noindent
    {\large\textbf{Estimators}}

    \textit{Estimator:} func $ T = g(X_1, \ldots, X_n) $ used to estimate 
    population param. $\theta$

    \textit{Estimate:} particular value the estimator takes for a particular
    realised sample by applying the function to dataset $t = g(x_1, \ldots,
    x_n)$. typically notated as a $t$ or by $\hat{\theta}$

    \textit{Unbiased} $T_n$ of $\theta$ if at any sample size $n$, $\E[T_n] = \theta$.
    Asymptotically unbiased if $\lim_{n \to \infty} \E[T_n] = \theta$.

    $T_n$ is a consistent estimator of $\theta$, if 
    \setlist{nolistsep}
    \begin{enumerate}[noitemsep]
        \item $T_n$ is asymptotically unbiased
        \item $\Var(T_n) \to 0$ as $n \to \infty$. 
    \end{enumerate}

    Suppose $X_1 , \ldots , X_n$ represent the random variables of a sample
    from a population. The sample mean estimator is defined by $\bar{X}$, where,
    \[
        \bar{X} = \frac{1}{n} \sum_{i=1}^n X_i.
    \] $\E[\bar{X}] = \mu, \Var(\bar{X}) = \frac{\sigma^2}{n}.$

    The sample variance estimator is denoted by $S^2$ and given by, \[
        S^2 = \frac{1}{n-1} \sum_{i=1}^n(X_i-\bar{X})^2.
    \] $\E[S^2] = \sigma^2.$

    $\bar{X}$ is an unbiased and consistent estimator for $\mu$. 
    $S^2$ is an unbiased estimator for $\sigma^2$. Also consistent if all fourth
    moments are equal and finite. 

    Standard Error $\text{S.E.}(\bar{X}) = \sqrt{\Var(\bar{X})} = \frac{\sigma}{\sqrt{n}}.$

    \noindent
    \textbf{Constructing an estimator}

    \textbf{Method of moments}

    Consider some (population) random variable $X$ with probability density/mass
    function $f_X(x|\theta)$ that depends on some parameter $\theta$. The $k$th
    population
    moment, for $k \in \N$, is defined as the expectation:
    \begin{equation*}
        \mu_k := \E [X^k] = \begin{cases}
            \sum_x x^k f_X(x | \theta), & \text{DRV $X$}, \\
            \int_{-\infty}^\infty x^k f_X(x | \theta) dx, & \text{CRV $X$.}
        \end{cases}
    \end{equation*}
    Let the random variables $X_1 , \ldots , X_n$ denote a potential sample from the
    population, then the $k$th sample moment estimator is: \[
        \overline{X^k} = \frac{1}{n} \sum^n_{i=1} X^k_i.
    \]
    Applying this transformation to a given sample, $x_1 , \ldots , x_n$, obtains the
    $k$th sample moment estimate $\bar{x^k}$.

    Let $x_1 , \ldots , x_n$ be a random sample from a distribution depending on
    the population parameters $\theta_1 , \ldots , \theta_n$ whose values are
    unknown. Suppose
    that these parameters can be expressed as some function of the first $m$ 
    population moments: \[ 
        (\theta_1 , \ldots , \theta_n) = h(\mu_1 , \ldots , \mu_m).
    \]

    The method of moments estimators $T_1, \ldots, T_m$ for the parameters 
    $\theta_1, \ldots, \theta_m$ are then constructed by applying the
    transformation $h$ with the corresponding sample moments: \[
        (T_1, \ldots, T_m) = h(\bar{X}, \ldots, \overline{X^m}). 
    \]

    \textbf{Advantages} Fairly simple and requires few assumptions about the
    pop dist; also consitent under a set of rel weak assumptions. 

    \textbf{Limitations} Can only be applied if an expressions for pop moments
    exists. Typically biased for a limited sample size. May not always satisfy
    any constraint on the parameter and sample size. 

    \textbf{Maximum likelihood}

    The \textit{likelihood function} is mathematically identical to the joint
    probability mass/density function for a sample $x_1 , \ldots , x_n$, but
    instead viewed as a function of the population parameter $\theta$. For an
    independent and identically distributed sample, the likelihood function is:
    \[
        L(\boldsymbol{\theta} | \mathbf{x}) = \prod_{i=1}^n f(x_i | \boldsymbol{\theta}).
    \]

    The log-likelihood function is simply the natural logarithm of the
    likelihood function:
    \begin{align*}
        \ell(\boldsymbol{\theta} | \mathbf{x}) &= \log{L(\boldsymbol{\theta} | \mathbf{x})} \\
                                               &= \log\left(\prod_{i=1}^n f(x_i | \boldsymbol{\theta}) \right) \\
                                               &= \sum_{i=1}^n \log{f(x_i | \boldsymbol{\theta})}
    \end{align*}

    The \textit{maximum likelihood estimate} for the population parameter
    $\theta$ is the function of the sample $x_1 , \ldots , x_n$ that maximises
    the likelihood function:
    \[
        \hat{\theta} = \arg\max_\theta L(\theta | \mathbf{x}).
    \] In practice, find the derivative of $\ell$ with respect to the $\theta$, 
    then set it to zero and work out what $\theta$ is. Confirm that it actually 
    maximises the function (2nd derivative is less than zero with $\theta$
    subbed in).

    \textbf{Advantages} Calculus can be applied to vast majority of likelihood
    functions. Even if differentiating is difficult, then it can be computed
    using numerical optimisation algorithms. Although often biased, they are all 
    consistent. Also invariant to re-parameterisation. 

    \textbf{Limitations} Requires full definition of the joint probability
    distribution for the sample. Needs generally pretty strong assumptions
    (e.g. indep and follow a stated distribution)

    \noindent
    \textbf{Interval Estimation}

    Let $\mathbf{X} = (X_1 , \ldots, X_n)$ represent the vector of sample random
    variables from some population with unknown parameter $\theta$. For some
    $\alpha \in [0, 1]$ (typically $\alpha = 0.05$), a $100(1 - \alpha)\%$
    confidence interval for the
    parameter $\theta$ is the interval with lower and upper bound estimators 
    $l(\mathbf{X})$ and $u(\mathbf{X})$ such that: \[
        \pf(l (\mathbf{X}) < \theta < u(\mathbf{X})) = 1 - \alpha. 
    \] It then follows that
    the $100(1 - \alpha)\%$ confidence interval for $\theta$ from the dataset 
    $x = (x_1 , \ldots , x_n )$ is computed as the interval $(l (\mathbf{x}),
    u(\mathbf{x}))$. 
    Generally will only consider the intervals where remaining percentage is split 
    equally above and below, i.e. 
    $\pf(l(\mathbf{X} > \theta)) = \pf(u(\mathbf{X} < \theta)) = \frac{\alpha}{2}$

    What is important to note is that this interval contains the unknown fixed parameter
    with a probability $95\%$ under the principle of repeated sampling. 

    \noindent
    \textbf{Normal dist with known $\sigma^2$}

    Let $Z = \frac{\bar{X} - \mu}{\sqrt{\sigma^2/n}}$ where $Z \sim \text{N}(0,1)$. 
    Then the $(1 - \alpha)\%$ confidence interval for $\mu$ is given by: 
    \[
        \left( \bar{x} - z_{\alpha/2} \frac{\sigma}{\sqrt{n}}, \bar{x} + z_{\alpha/2} \frac{\sigma}{\sqrt{n}} \right).
    \]

    Use \verb|qnorm(p, mean, sd)| to calculate the quantiles for this. 
    For 95\% interval, let $p = 0.975$.

    \noindent
    \textbf{Confidence interval via the central limit theorem}

    As long as $X_1, \ldots, X_n$ are iid, with known $\sigma^2$, then eventually 
    we have that $\bar{X}$ converges toward std norm dist according to CLT, i.e.
    \[
        \frac{\bar{X} - \mu}{\sqrt{\Var(\bar{X})}} \sim \text{N}(0,1) \text{ as } n \to \infty.
    \]

    Given $n$ sufficiently large, we have that the confidence interval is \[
        \left( \bar{X} - z_{\alpha / 2} \frac{\sigma}{\sqrt{n}},\bar{X} + z_{\alpha / 2} \frac{\sigma}{\sqrt{n}} \right)
    \]

    If we don't know what $\sigma^2$ is then we can use $S^2$ to estimate it. In this case: 
    \[
        \left( \bar{x} - z_{\alpha / 2} \sqrt{\frac{s^2}{n}},\bar{x} + z_{\alpha / 2} \sqrt{\frac{s^2}{n}} \right)
    \]

    \noindent
    \textbf{Confidence interval for sample mean with small sample size}

    Let $Z_1, \ldots, Z_n$ be indep std norm RVs. Define $X = \sum^n_{i=1} Z^2_i$, 
    then the RV $X$ has a chi-squared distribution with $n$ degrees of freedom, $X \sim \chi^2_n$.

    Let $Z = \text{N}$ and $Y \sim \chi^2_m$ be indep RVss, then Student's $t$-dist
    on $m$ degrees of freedom is: \[
        \frac{Z}{\sqrt{Y/m}} \sim t_m.
    \] It is symmetric at $0$, more mass in tails compared to std norm. 

    Let $X_1, \ldots, X_n$  be indep normal RVs, with unknown $\mu, \sigma^2$. Then: 
    \begin{gather*}
        \bar{X} \& S^2 \text{ are independent} \\
        \frac{\bar{X} - \mu}{\sqrt{\sigma^2 / n}} \sim \text{N}(0,1) \\
        \frac{(n-1)S^2}{\sigma^2} \sim \chi^2_{n-1}
    \end{gather*}
    Therefore we can define RV $T$, indep of $\sigma$, that follows Student's $t$-dist
    on $n-1$ degrees of freedom: \[
        T = \frac{\bar{X} - \mu}{\sqrt{S^2/n}} \sim t_{n-1}.
    \]

    Confidence interval is: \[
        \left( \bar{x} - t_{n-1 ; \frac{\alpha}{2}}\sqrt{\frac{s^2}{n}}, \bar{x} + t_{n-1 ; \frac{\alpha}{2}}\sqrt{\frac{s^2}{n}}\right)
    \]
    
    \noindent
    \textbf{Confidence interval for $S^2$}

    Since, $Y = \frac{(n-1)S^2}{\sigma^2} \sim \chi^2_{n-1}$, we need to find 
    $\pf\left( \chi^2_{n-1; 1- \frac{\alpha}{2}} \leq Y \leq \chi^2_{n-1; \frac{\alpha}{2}} \right)$

    So for a given dataset the confidence interval is: 
    \[ 
    \left( \frac{(n-1)s^2}{\chi^2_{n-1; \frac{\alpha}{2}}}, \frac{(n-1)s^2}{\chi^2_{n-1; \frac{\alpha}{2}}}\right)
    \]

% Hypothesis Testing

    \noindent
    {\large\textbf{Hypothesis testing}}

    Steps: 
    \setlist{nolistsep}
    \begin{enumerate}[noitemsep]
        \item State the null and alt hypothesis.
        \item Select appropriate test statisticand determine dist assuming $H_0$ is true. Under this set the decision rule.
        \item Compute the observed test statistic from the sampled data. 
        \item Compare $t$ against decision rule and give a concluding statement.
    \end{enumerate}

    One sided alt asserts that the deviation from the null occurs in a specific direction. 
    Two sided alt specifies that the parameter may differ in any direction. 

    Can only either reject the null hypothesis or fail to reject the null hypothesis. 

    A type I error (or false-positive error) occurs when the null hypothesis is
    incorrectly rejected when it is actually true.

    The probability of incurring a Type I error, 
    $\pf(\text{Reject }H_0 |H_0
    \text{is true})$, is called the significance level of the hypothesis test,
    and is commonly denoted by $\alpha$.

    A type II error, or false-negative, occurs when we incorrectly fail to
    reject the null hypothesis when the null hypothesis is actually incorrect.

    The probability of incurring a Type II, $$\pf(\text{Fail to reject } H_0
    |H_0 \text{ is false})$$ is typically denoted by $\beta$.

    The complement of this probability $$\pf(\text{Reject } H_0 |H_0 \text{ is false}) = 1 - \beta,$$
    is called the power of the test.

    The $p$-value is the probability of obtaining a value for the test statistic
    random variable, $T$, that is as or more extreme than the observed test
    statistic, $t$, under the assumption of repeated sampling under the null
    hypothesis.

    The effect of the phrase `as or more extreme than' on computing the p-value
    depends on the type and direction of the alternative hypothesis:
    \setlist{nolistsep}
    \begin{itemize}[noitemsep]
        \item one-sided in the positive direction, $H_1 : \theta > \theta_0 \implies p = \pf (T \geq t | T \sim H_0 )$;
        \item one-sided in the negative direction, $H_1 : \theta < \theta_0 \implies p = \pf (T \leq t | T \sim H_0 )$;
        \item Two sided, $H_1 : \theta \neq \theta_0 \implies p = 2 \min \{\pf (T \leq t | T \sim H_0 ) , \pf (T \geq t | T \sim H_0 )\}$.
    \end{itemize}

    \noindent
    \textbf{Power of a hypothesis test} 

    Increasing significance level, sample size, and effect size increases the power of a hypothesis test. 

    \noindent 
    \textbf{$Z$-test}

    $X_1, \ldots, X_n$  are indep $\text{N}(\mu, \sigma^2)$ RVs where $\sigma^2$ is known. Use this test statistic: 
    \[
        Z = \frac{\bar{X} - \mu_0}{\sqrt{\sigma^2/n}} \sim \text{N}(0, 1),
    \] where we have $H_0: \mu = \mu_0, H_1: \mu \neq \mu_0.$
    Then the critical region is $C = \{ z : |z| \geq z_0\}$, $\pf(Z \geq z_0) = \frac{\alpha}{2}$.

    Seems that these are generally two sided. If $z \in C$ then reject. 

    \textbf{Example:} $X_1, \ldots, X_100$, $X_i \sim \text{N}(\mu, 44)$ iid, 
    then $\bar{X} \sim \text{N}(\mu, 1.44).$ $H_0: \mu = 65$, $H_1: \mu \neq 65.$
    $Z = \frac{\bar{X} - 65}{1.2} \sim \text{N}(0,1).$ At $5\%$ sig level, we 
    have that $C = \{z : |z| \geq 1.96\}$. Calculated $z = -1.25$, $z \notin C$, so 
    fail to reject. Or, with $p$-value, we have 
    \begin{align*}
        \pf(|Z| \geq |-1.25|) &= 2\pf(Z \leq -1.25) \\
                              &= 2\pf(Z \geq 1.25) \\
                              &\approx 0.2113. > 0.05.
    \end{align*}

    \noindent
    \textbf{One-sample-$t$-test}

    $X_1, \ldots, X_n$ be indep $N(\mu, \sigma^2)$ RV, with $\sigma^2$ unknown. 
    Test statistic: $T = \frac{\bar{X} - \mu_0}{\sqrt{S^2 / n}}$.
    For tails: $\pf( T \geq t_0 | T \sim t_{n-1}) = \frac{\alpha}{2}.$

    \noindent
    \textbf{Paired $t$-test}

    Like the one sample test but we have $D_i = X_i - Y_i$ where $\mu_D = \mu_X - \mu_y$. 
    $\sigma^2_D$ is the unknown variance of the differences. 
    
    Test statistic: $T = \frac{\bar{D} - \mu_0}{\sqrt{S^2 / n}} \sim t_{n-1}$.

    \noindent
    \textbf{two sample $t$ test} 

    $X_1, \ldots, X_n$, $Y_1, \ldots, Y_m$  
    We have $H_0: \mu_X - \mu_Y = \Delta_0, H_1: \mu_X - \mu_Y \neq \Delta_0$.
    So 
    \begin{gather*} 
        \bar{X} \sim N(\mu_X, \frac{\sigma^2}{n}, \bar{Y} \sim N(\mu_Y, \frac{\sigma^2}{m} \\
        \bar{X} - \bar{Y} \sim N(\mu_X - \mu_Y, \sigma^2(\frac{1}{n} + \frac{1}{m})) \\
        Z = \frac{(\bar{X} - \bar{Y}) - \Delta_0}{\sqrt{\sigma^2 (\frac{1}{m} + \frac{1}{n}}} \sim N(0, 1). \\
        S^2_p = \frac{\sum_{i=1}^{n}(X_i - \bar{X})^2 + \sum_{j=1}^m(Y_j - \bar{Y})^2}{m+n-2} \\
        S^2_p = \frac{(n-1)S^2_X + (m-1)S^2_Y}{m + n - 2}
    \end{gather*}

    $S_P^2$ is an unbiased estimator for $\sigma^2$

    Test statistic: $T = \frac{\bar{X} - \bar{Y} - \Delta_0}{\sqrt{S^2_P(\frac{1}{m} + \frac{1}{n}}} \sim t_{m+n-2} $

    \noindent
    \textbf{$F$-test for equal variances}

    Assessment of whether the two population variances are equal. Wish to test: 
    \[ H_0: \sigma^2_X = \sigma^2_Y, H_1: \sigma^2_X \neq \sigma^2_Y.\]

    Let $U \sim \chi^2_m$, and $V \sim \chi^2_n.$ Then 
    \[ F = \frac{U / m}{V / n} \sim F_{m,n}.\]

    Test statistic: $F = \frac{S^2_X}{S^2_Y} \sim F_{n-1, m-1}.$ 

    Would need to do \verb|qf(c(lower, upper), df1, df2)|. 

    \noindent
    \textbf{Assessing normality}

    \textbf{The empirical cumulative distribution function}

    Consider the observed sample $x_1, \ldots, x_n$ from some common random
    variable $X$. Let the notation $x_{(i )}$ denote the $i$th sample in ascending
    order so that $x_{(i )} \leq x_{(i +1)}$ with $x_{(1)}$ as the minimum value and $x_{(n)}$ 
    is the largest. The empirical cumulative distribution function is a
    step-wise function mapping the sample space to the unit space of cumulative
    probabilities, $\hat{F}_X : \R \to [0, 1]$. Specifically, it is defined as: \[
    \hat{F}_X(x) = \frac{1}{n} \sum^n_{i=1} \mathbb{I} [ x_{(i)} \leq x]. \]

    From this we see that the samples should follow a roughly `S' shaped curve.  

    Generated using \verb|ecdf()|.

    For quantile-quantile plot do \verb|qqnorm(x)| and \verb|qqline(x)|.

    Dots should be roughly along a straight line.

    \noindent
    {\large\textbf{Linear regression}}

    Consider the simple linear regression model with a single explanatory
    variable, $\E [Y ] = \alpha+\beta x$. Suppose we have a data set consisting
    of the
    explanatory and response data pairs $(x_1 , y_1 ), \ldots , (x_n , y_n )$,
    then the least-squares estimates, $\hat{\alpha}$ and $\hat{\beta},$ that
    minimise the sum
    of squares
    formula are:
    \begin{gather*}
        \hat{\alpha} = \bar{y} - \hat{\beta}\bar{x} \\
        \hat{\beta} = \frac{\sum_{i=1}^n (x_i - \bar{x})(y_i - \bar{y})}{\sum_{i=1}^n (x_i - \bar{x})^2} = \frac{s_{xy}}{s^2_x} = r\frac{s_{y}}{s_x} \\
        Q(\alpha,\beta) = \sum_{i=1}^n \{y_i - (\alpha + \beta x_i)\}^2
    \end{gather*}
    The overall discrepancy between the data and the model can be expressed by the
    sum of squares formula, $Q(\alpha, \beta)$. 

    Consider the simple linear regression model $Y_i = \alpha + \beta x_i + \epsilon_i$ for 
    $i = 1, \ldots , n$ with independent error random variables that have zero expectation,
    $E [\epsilon_i ] = 0$, and constant variance, $\Var (\epsilon_i ) =
    \sigma^2 $. Both least-squares estimators are unbiased and consistent in
    estimating the regression parameters.

    The fitted values, or predicted values, $\hat{Y}_1 , \ldots , \hat{y}_n$
    are the response
    values obtained by substituting the corresponding explanatory values into
    the fitted regression line: 
    \[ \hat{y}_i = \hat{\alpha} + \hat{\beta}x_i \text{ for } i = 1, \ldots, n. \]

    The residuals, $e_1 , \ldots , e_n$ , are the differences between the observed
    and predicted response values:
    \[ e_i = y_i - \hat{y}_i = y_i - \hat{\alpha} - \hat{\beta}x_i \text{ for } i = 1, \ldots, n.\]

    The residual sum of squares is defined to be the sum of the squared errors
    between the observed responses $y_i$ and the predicted value $\hat{y}_i$ from the
    simple linear regression model, $\sum_{i=1}^n (y_i - \hat{y}_i )^2$. From
    this, we define
    the least-squares estimate of $\sigma^2$ , also known as the residual variance, as:
    \[ s_e^2 = \frac{1}{n-2} \sum_{i=1}^n (y_i - \hat{y}_i)^2 \]
    This is unbiased in estimating $\sigma^2$. 
    \begin{gather*}
        \widehat{\text{S.E.}}(\hat{\alpha}) = s_e\sqrt{\frac{1}{n} + \frac{\bar{x}^2}{(n-1)s_x^2}}, \qquad \\
        \widehat{\text{S.E.}}(\hat{\beta}) = \frac{s_e}{\sqrt{(n-1)s_x^2}}
    \end{gather*}

    \begin{gather*}
        Y_i \overset{\text{indep.}}{\sim} N(\alpha + \beta x_i,\, \sigma^2), \quad i = 1, \ldots, n \\
        \hat{\alpha} \sim N\!\left(\alpha,\; \sigma^2\!\left[\frac{1}{n} + \frac{\bar{x}^2}{(n-1)s_x^2}\right]\right), \qquad \\
        \hat{\beta} \sim N\!\left(\beta,\; \frac{\sigma^2}{(n-1)s_x^2}\right) \\
    \end{gather*}

    The maximum likelihood estimates for the regression parameters $\alpha$ and $\beta$ in
    the normal linear regression model are the least-squares estimates $\hat{\alpha}$ and
    $\hat{\beta}.$ 

    The distribution of residual variance estimator is: 
    \[ \frac{(n-2)S_e^2}{\sigma^2} \sim \chi^2_{n-2}. \]
 
    $\hat{\alpha}$ and $\hat{\beta}$ are independent of $S^2_e$, marginal dist: 
    \begin{gather*}
        \frac{\hat{\alpha} - \alpha}{\sqrt{S_e^2\!\left(\frac{1}{n} + \frac{\bar{x}^2}{(n-1)s_x^2}\right)}} \sim t_{n-2}, \qquad \\
        \frac{\hat{\beta} - \beta}{\sqrt{\frac{S_e^2}{(n-1)s_x^2}}} \sim t_{n-2} 
    \end{gather*}

    The 95\% confidence interval for $\alpha$ and $\beta$ are: 
    \begin{gather*}
        \hat{\alpha} \pm t_{n-2;\,0.025}\,s_e\sqrt{\frac{1}{n} + \frac{\bar{x}^2}{(n-1)s_x^2}}, \\
        \hat{\beta} \pm t_{n-2;\,0.025}\,\frac{s_e}{\sqrt{(n-1)s_x^2}}
    \end{gather*}

    \begin{verbatim}
Usage: lm(formaula, data)
Inputs: 
formula - the linear model expression of the 
form y ~ x
data - a data frame containing the variables 
used in the formula
    \end{verbatim}
    Use \verb|summary(model)| in order to get information about the 
    residuals, coefficients (least squares estimate for each regression
    param and their std error), residual std error, multiple $R^2$, and the 
    $F$-statistic (test stat and df for corresponding $F$-dist)
    
    In the case of a single explanatory var, you have $R^2 = \frac{s^2_{xy}}{s^2_x s^2_y}$
    Between $[0, 1]$. Want it to be close to 1. Should be used to compare, not say 
    if any one model is good or bad. 

    \noindent
    \textbf{Confidence interval for $\E[Y_0]$}

    Essentially for a given $x_0$ in the dataset already. 
    \[ \frac{\hat{\E}[Y_0] - \E[Y_0]}{\sqrt{S^2_e \left( \frac{1}{n} + \frac{(x_0 - \bar{x})^2}{(n-1)s_x^2}\right)}} \sim t_{n-2} \]
    \[ \hat{\alpha} + \hat{\beta}x_0 \;\pm\; t_{n-2;\,0.025}\,s_e\sqrt{\frac{1}{n} + \frac{(x_0-\bar{x})^2}{(n-1)s_x^2}} \]

    \noindent
    \textbf{Prediction interval for $\E[Y_0]$}

    For a potential future observation of $Y_0$ at $x_0$. 
    
    \begin{gather*}
        Y_0 = \alpha + \beta x_0 + \epsilon_0 \\
        \hat{Y}_0 = \hat{\alpha} + \hat{\beta}x_0 + \epsilon_0 = \hat{\mathbb{E}}[Y_0] + \epsilon_0 \\
        \text{Var}(\hat{Y}_0) = \sigma^2\!\left(1 + \frac{1}{n} + \frac{(x_0-\bar{x})^2}{(n-1)s_x^2}\right) \\
        \hat{\alpha} + \hat{\beta}x_0 \;\pm\; t_{n-2;\,0.025}\,s_e\sqrt{1 + \frac{1}{n} + \frac{(x_0-\bar{x})^2}{(n-1)s_x^2}}
    \end{gather*}

    \begin{verbatim}
Usage: 
predict(object (model that you already have), 
newdata, 
interval = c("none", "confidence", "prediction"),
level = 0.95)
Inputs: 
object - the linear regression object from lm()
newdata - a data frame with the new explanatory 
variables x_0
interval - type of interval to compute
level - confidence level of the interval
    \end{verbatim}
    This gives you the value(s) of $y_0$ for a given $x_0$

    \noindent 
    \textbf{Checking assumptions}

    The normal linear regression model specifies that $Y_i sim N (\alpha + \beta x_i , \sigma^2 )$ 
    for $i = 1, \ldots , n$. The assumptions made in defining this statistical
    model are:
    \begin{enumerate}
        \item Independence between response variables;
        \item Linearity of the expectation, $\E(Y_i ) = \alpha + \beta x_i$;
        \item Homoscedasticity (i.e. have the same variance); and,
        \item Normally distributed.
    \end{enumerate}

    \textbf{Residuals vs Fitted}
    This is a scatter plot of the residuals $e_1 , \ldots , e_n$ against fitted
    values $\hat{y}_1 , \ldots , \hat{y}_n$ . If the linear assumption $\E [Y ] = \alpha +
    \beta x$ is true
    then the residuals should be randomly scattered around zero with no
    discernible clustering or pattern. Furthermore, we can also use this plot
    to check the homoscedastic (constant variance) assumption to see whether
    the range of the scatter of points is consistent over the range of fitted
    values. This visual assessment can be challenging to perform accurately, so
    we may instead examine the next plot.

    \textbf{Scale-Location}
    This is a scatter plot of the transformation of the standardized residuals
    $\sqrt{|e_1^*|}, \ldots , \sqrt{|e_n^*|} $ against fitted values $\hat{y}_1 ,
    \ldots , \hat{y}_n$ . This is
    used to assess the homoscedastic (constant variance) assumption where, if
    true, we should expect the points to be randomly scattered around a
    constant value $(\sim 0.82)$ across the range of fitted values.

    \textbf{Normal QQ-plot}
    This is a normal QQ-plot of the standardised residuals$e_1^* , \ldots , e_n^*$ 
    where $e_i* = \frac{e_i}{s_e}$ for $i = 1, \ldots , n$. If the assumption that the
    response random variables are independent and normally distributed then we
    can consider the standardised residuals as a sample of $n$ values from the
    standard normal distribution.

    Can get these in \verb|R| using \verb|plot(model)|.

    Remember that this is correlation not causation. 

    \noindent 
    {\large\textbf{Multiple regression}}

    Consider the linear regression model $\underline{Y} = \mathbf{X}
    \underline{\beta}+ \underline{\epsilon}$, expressed in matrix notation,
    where $\underline{Y}$ is the response vector random variable, $\mathbf{X}$
    is the $n \times p$ design
    matrix containing the explanatory variables,
    $\underline{\beta}$ is the vector or regression parameters and
    $\underline{\epsilon}$ is the vector of
    random errors where each element is
    independent with $\E [\epsilon_i ] = 0$ and $\Var (\epsilon_i ) = \sigma^2$
    . The least-squares vector estimate $\hat{\beta} = (\hat{\beta}_1 , \ldots
    , \hat{\beta}_p )^T$ that
    minimises the sum of squared formula $Q(\underline{\beta}) = (\underline{y}
    - \mathbf{X} \underline{\beta})^T (\underline{y} - \mathbf{X} \underline{\beta})$ is:
    \[ \hat{\beta} = (\mathbf{X}^T \mathbf{X} )^{-1} \mathbf{X}^T \underline{y}. \]

    Need to ensure that all of the columns of $\mathbf{X}$ are linearly
    independent.

    Residual variance for multiple regression:
    \begin{align*}
    % Residual variance (multiple regression) \\
        S_e^2 &= \frac{1}{n-p} \sum_{i=1}^n \left[Y_i - (\hat{\beta}_1 x_{i,1} + \cdots + \hat{\beta}_p x_{i,p})\right]^2 \\
        &= \frac{1}{n-p}(\mathbf{y} - \mathbf{X}\hat{\boldsymbol{\beta}})^\top(\mathbf{y} - \mathbf{X}\hat{\boldsymbol{\beta}}) \\
    \end{align*}

    Marginal distribution of $\hat{\beta}_h$
    % Marginal distribution of beta_j-hat \\
    \[ \frac{\hat{\beta}_j - \beta_j}{\text{S.E.}(\hat{\beta}_j)} \sim t_{n-p} \] 
    Confidence interval for each $\hat{\beta}_h$.
    % 95% CI for each regression parameter \\
    \[ \hat{\beta}_j \pm t_{n-p;\,0.025}\;\widehat{\text{S.E.}}(\hat{\beta}_j), \quad j = 1, \ldots, p \]

    Adjusted $R^2$ takes into account the number of variables in regression model:
    % Adjusted R^2 \\
    \[ R^2_{\text{adj.}} = 1 - (1 - R^2)\frac{n-1}{n-p} \]

    \noindent
    \textbf{ANOVA $F$-test for model comparison}

    \begin{gather*}
        \text{ANOVA F-test definition:} \\
        H_0: \beta_{p_0+1} = \cdots = \beta_{p_1} = 0, \quad \\
        H_1: \exists\, j \in \{p_0+1,\ldots,p_1\} \text{ s.t. } \beta_j \neq 0 \\
        F = \frac{(\text{RSS}_0 - \text{RSS}_1)\,/\,(p_1 - p_0)}{\text{RSS}_1\,/\,(n - p_1)} \sim F_{p_1-p_0,\,n-p_1} \\
        \text{RSS}_1 = (n - p_1)S_1^2, \qquad \text{RSS}_0 = (n - p_0)S_0^2 \\
        \text{Categorical binary indicator:} \\
        a_{i,j} = \begin{cases} 1, & \text{if sample } i \text{ belongs to category } j \\ 0, & \text{otherwise} \end{cases} \\
        \text{One-way ANOVA model:} \\
        Y_{i,j} \overset{\text{iid}}{\sim} N(\mu_j,\, \sigma^2), \quad i = 1,\ldots,n_j,\; j = 1,\ldots,k \\
        \hat{\mu}_j = \bar{y}_j = \frac{1}{n_j}\sum_{i=1}^{n_j} y_{i,j} \\
        s_e^2 = \frac{1}{n-k}\sum_{j=1}^k \sum_{i=1}^{n_j}(y_{i,j} - \bar{y}_j)^2 \\
        \text{ANOVA decomposition (one-way):} \\
        \sum_{j=1}^k \sum_{i=1}^{n_j}(Y_{i,j} - \bar{Y})^2 \\
        = \sum_{j=1}^k n_j(\bar{Y}_j - \bar{Y})^2 \\
        + \sum_{j=1}^k \sum_{i=1}^{n_j}(Y_{i,j} - \bar{Y}_j)^2 \\
        \text{RSS}_y = \text{RSS}_b + \text{RSS}_e \\
        (n-1)S_y^2 = (k-1)S_b^2 + (n-k)S_e^2 \\
        F = \frac{\text{RSS}_b\,/\,(k-1)}{\text{RSS}_e\,/\,(n-k)} = \frac{S_b^2}{S_e^2} \sim F_{k-1;\,n-k} \quad \text{under } H_0 \\
    \end{gather*}
    \begin{gather*}
        \text{Two-way ANOVA model:} \\
        \mu_{i,j} = \mu + \alpha_i + \beta_j \\
        \sum_{i=1}^b \alpha_i = 0, \qquad \sum_{j=1}^k \beta_j = 0 \\
        \hat{\mu} = \bar{y}, \quad \hat{\alpha}_i = \bar{y}_{i,\bullet} - \bar{y}, \quad \hat{\beta}_j = \bar{y}_{\bullet,j} - \bar{y} \\
        \bar{y}_{i,\bullet} = \frac{1}{k}\sum_{j=1}^k y_{i,j}, \qquad \\
        \bar{y}_{\bullet,j} = \frac{1}{b}\sum_{i=1}^b y_{i,j}, \qquad \\
        \bar{y} = \frac{1}{n}\sum_{i=1}^b\sum_{j=1}^k y_{i,j} \\
        \text{Two-way ANOVA decomposition:} \\
        \sum_{i=1}^b\sum_{j=1}^k (Y_{i,j} - \bar{Y})^2 \\
        = k\sum_{i=1}^b(\bar{Y}_{i,\bullet} - \bar{Y})^2 \\
        + b\sum_{j=1}^k(\bar{Y}_{\bullet,j} - \bar{Y})^2 \\
        + \sum_{i=1}^b\sum_{j=1}^k (Y_{i,j} - \bar{Y}_{i,\bullet} - \bar{Y}_{\bullet,j} + \bar{Y})^2 \\
        \text{RSS}_y = \text{RSS}_b + \text{RSS}_t + \text{RSS}_e \\
        (n-1)S_y^2 = (b-1)S_b^2 + (k-1)S_t^2 + (n-b-k+1)S_e^2 \\
        F_b = \frac{\text{RSS}_b\,/\,(b-1)}{\text{RSS}_e\,/\,(n-b-k+1)} = \frac{S_b^2}{S_e^2} \sim F_{b-1;\,n-k-b+1} \\
        F_t = \frac{\text{RSS}_t\,/\,(k-1)}{\text{RSS}_e\,/\,(n-b-k+1)} = \frac{S_t^2}{S_e^2} \sim F_{k-1;\,n-k-b+1}
    \end{gather*}

% Common Distributions

    \noindent
    {\large\textbf{Common Distributions}}

    \noindent
    \textbf{Discrete}

    \textbf{Bernoulli}, $p \in [0, 1]X \sim \text{Bern}(p).$

    Params: Probablility, .
    $f_X(x) = p^x(1-p)^{1-x}$, for $x \in \{0,1\}$.
    $\E[X] = p, \Var(X) = p(1-p)$.

    \textbf{Binomial}, $X \sim \text{Binom}(n, p)$.

    Params: Sample Size, $n \in \N$; Probablility, $p \in [0, 1]$.
    $f_X(x) = \binom{n}{x}p^x(1-p)^{n-x}$, for $x \in \{0, 1, \ldots, n\}$.
    $\E[X] = np, \Var(X) = np(1-p)$.

    \textbf{Geometric}, $X \sim \text{Geom}(p)$.

    Params: Probablility, $p \in [0, 1]$.
    $f_X(x) = p(1-p)^{x - 1}$, for $x \in \{1, 2, \ldots\}$.
    $\E[X] = \frac{1}{p}, \Var(X) = \frac{1-p}{p}$.

    \textbf{Poisson}, $X \sim \text{Pois}(\lambda)$.
    
    Params: Rate, $\lambda > 0$.
    $f_X(x) = \frac{\lambda^x \exp(-\lambda)}{x!}$, for $x \in \{0, 1, 2, \ldots\}$.
    $\E[X] = \lambda, \Var(X) = \lambda$.

    \noindent
    \textbf{Continuous}

    \textbf{Uniform}, $X \sim \text{Unif}(a, b)$.

    Params: Limits, $a, b \in \R$.
    $f_X(x) = \frac{1}{b - a}$, for $x \in [a, b]$.
    $\E[X] = \frac{a + b}{2}, \Var(X) = \frac{(b - a)^2}{12}$.

    \textbf{Exponential}, $X \sim \text{Exp}(\lambda)$.

    Params: Rate, $\lambda > 0$.
    $f_X(x) = \lambda \exp(-\lambda x)$, for $x \in [0, \infty)$.
    $\E[X] = \frac{1}{\lambda}, \Var(X) = \frac{1}{\lambda^2}$.

    \textbf{Gamma}, $X \sim \Gamma(\alpha, \beta)$.

    Params: Shape, $\alpha > 0$; Rate, $\beta > 0$.
    $f_X(x) = \frac{\beta^\alpha}{\Gamma(\alpha)} x^{\alpha - 1} \exp(-\beta x)$, for $x \in [0, \infty)$.
    $\E[X] = \frac{\alpha}{\beta}, \Var(X) = \frac{\alpha}{\beta^2}$.

    \textbf{Beta}, $X \sim \text{Beta}(\alpha, \beta)$.

    Params: Shapes, $\alpha, \beta > 0$.
    $f_X(x) = \frac{x^{\alpha - 1}(1-x)^{\beta - 1}}{B(\alpha, \beta)}$, 
    for $x \in [0, 1]$.
    $\E[X] = \frac{\alpha}{\alpha + \beta}, 
    \Var(X) = \frac{\alpha\beta}{(\alpha + \beta)^2(\alpha + \beta + 1)}$.

    \textbf{Normal}, $X \sim \text{N}(\mu, \sigma^2)$.

    Params: Expectation, $\mu \in \R$; Std Deviation, $\sigma > 0$.
    $f_X(x) = \frac{1}{\sqrt{2 \pi \sigma^2}}\exp\left\{ \frac{-(x - \mu)^2}{2
    \sigma^2} \right\}$, for $x \in \R$.
    $\E[X] = \mu, \Var(X) = \sigma^2$.
    \textit{Sum of normals is normal, note that $\bar{X} \sim \text{N}(\mu, \frac{\sigma}{n}$}

    \textbf{Chi-Squared}, $X \sim \chi^2_{v}$.

    Params: Degree of Freedom, $v > 0$.
    $f_X(x) = \frac{1}{2^{\frac{v}{2}}\Gamma\left(\frac{k}{2}\right)}
            x^{\frac{v}{2}-1} \exp\left(\frac{x}{2}\right),$
    for $x \in [0, \infty)$.
    $\E[X] = v, \Var(X) = 2v$.

    \textbf{Student's $\mathbf{t}$}, $X \sim t_v$.

    Params: Degree of Freedom, $v > 0$.
    $f_X(x) = \frac{1}{\sqrt{v}B\left(\frac{v}{2},\frac{1}{2}\right)}
    \left(1 + \frac{x^2}{v}\right)^{-\frac{v+1}{2}},$
    for $x \in \R$.
    $\E[X] = 0, \Var(X) = \frac{v}{v-2}(v > 2)$.

    \textbf{Cauchy}, $X \sim \text{Cauchy} (X \sim t_1)$.

    No parameters. 
    $f_X(x) = \frac{1}{\pi (1 + x^2)}$ for $x \in \R$.
    Both $\E[X]$ and $\Var(X)$ are undefined.

    \textbf{F}, $X \sim \text{F}(\lambda, v)$.

    Params: Degrees of Freedom, $\lambda, v > 0$.
    $f_X(x) = \frac{\lambda^{\frac{\lambda}{2}}v^{\frac{v}{2}}}{B\left(\frac{\lambda}{2},\frac{v}{2}\right)}
    x^{\frac{\lambda}{2} - 1}(\lambda x + v)^{- \frac{\lambda + v}{2}}$,
    for $x \in [0, \infty)$. 
    $\E[X] = \frac{v}{v - 2} (v > 2), 
    \Var(X) = \frac{2v^2(\lambda + v - 2)}{\lambda(v - 2)^2(v - 4)} (v > 4) $

% Common R Commands

    \noindent
    {\large\textbf{Common R Commands}}

    \verb|qnorm(p, mean, sd)|
    \verb|qt(p, df)|

    \verb|q-| is the quantile function
    \verb|p-| is the cdf
    \verb|d-| gives the density
    \verb|r-| gives random values

    \begin{verbatim}
Usage: t.test(x, y = NULL, 
alternative = c("two.sided", "less", "greater"),
mu = 0, 
paired = FALSE, 
var.equal = FALSE, 
conf.level = 0.95)

Inputs: x, y - Vectors containing the sample data. 
alternative - Type of alternative hypothesis, 
two-sided stated by default. 
mu - The expected value assumed under the null 
hypothesis. 
paired - If TRUE, then performing a paired t-test. 
var.equal - If TRUE, then populations have equal 
variance in 2-sample t-test. 
conf.level - Confidence level of the test, equal 
to ‘1 - significance level’.
    \end{verbatim}
\end{multicols}

\end{document}
